\documentclass[11pt]{article}

\usepackage[margin=25mm]{geometry}

\usepackage{amsmath,amssymb,fontspec,unicode-math}

\IfFontExistsTF{DejaVu Serif}{\setmainfont{DejaVu Serif}}{\setmainfont{Latin Modern Roman}}

\IfFontExistsTF{Latin Modern Math}{\setmathfont{Latin Modern Math}}{\setmathfont{DejaVu Math TeX Gyre}}

\usepackage{xurl,hyperref}

\setlength{\parindent}{0pt}\setlength{\parskip}{6pt}

\begin{document}

\title{Visual Instruction Bottleneck Tuning}

\author{ELVIA preliminary AI analysis}\date{21 October 2026}\maketitle

Preliminary assessment: Major.\par

Independent scientific review remains pending. Author adoption has not been requested.

The actual interpolated decoder law is not controlled by the auxiliary KL, so the claimed common-latent information-bottleneck lower bound fails. Public KL/routing bugs and deletion of a valid bound’s residual term have separate scopes.

\section*{Review scope}

Latest-version, claim-focused review of the cited author-paper passages and applicable current author-release paths. The stated counterexamples and scope controls were checked analytically; no author program, old verifier, training run or benchmark was executed.

\section*{Problem 1: The released Gaussian KL halves the log-variance term}

Source: Appendix A.2, PDF p.28, Eq.(8); vittle\_llama.py, lines 252–289; repository issue \#1

Assessment: supported; impact: minor

\subsection*{Source claim (paraphrase)}

The paper’s Gaussian KL is correct, but `vib\_kld\_loss` and `vib\_kld\_loss\_lp` use half the required additive log-variance ratio.

\subsection*{Analysis}

At zero means, posterior variance \(1/2\), prior variance \(1\), the release returns \(-0.0767132\) instead of \(0.0965736\), within its clamps. At matching log variance \(s=0\), its derivative is \(1/4\) instead of zero; the fixed-prior minimum shifts to variance \(1/2\). The learned-prior mismatch was already raised in repository issue \#1 on 28 January 2026.

\subsection*{Finding and scope}

The later release computes neither the stated KL values nor gradients. Original checkpoint training is not identified.

\subsection*{Proposed correction}

Use \(\tfrac12[\log(v_p/v_q)-1+(v_q+(m_q-m_p)^2)/v_p]\), with consistent coordinate reduction.

Repository issue \#1 already identifies the learned-prior coefficient mismatch; no author adoption of this correction is established.

\section*{Problem 2: Sequence length does not identify modality positions}

Source: Section 3.3, PDF p.5; Appendix A.2, PDF pp.28–29; vittle\_arch.py packing and vittle\_llama.py, lines 312–334 and 439–490

Assessment: supported; impact: minor

\subsection*{Source claim (paraphrase)}

The public packing path replaces image sentinels in place, while the bottleneck route infers a visual prefix from total-length subtraction.

\subsection*{Analysis}

IDs \([1,2,\mathrm{image},3,4]\) expanded with visual rows \([101,102,103]\) become \([1,2,101,102,103,3,4]\). True visual indices are \(\{2,3,4\}\), but the inferred count \(7-5=2\) selects prefix \(\{0,1\}\), both text. Replacing one sentinel increases length by \(n-1\), not \(n\), even when the image is first. Embedding-only input sets the inferred visual count to zero in the inspected branch.

\subsection*{Finding and scope}

The public route does not implement the intended modality split. Full intended-version generation behavior and benchmark effects are unmeasured.

\subsection*{Proposed correction}

Carry an explicit visual-token mask through expansion, padding, truncation and cached generation; exclude padding separately.

An explicit local correction is described; author adoption and a corresponding published erratum are not established.

\section*{Problem 3: The skip contribution breaks the common-latent lower bound}

Source: Eqs.(2)–(5), PDF p.4; Section 3.3 and footnote 2, PDF p.5; Implication paragraph, PDF p.6; Appendix D, PDF p.38

Assessment: supported; impact: major

\subsection*{Source claim (paraphrase)}

Section 3.2 introduces Eq.(4) as a variational IB lower bound and Eq.(5) as its empirical estimate. The implemented likelihood receives \(\hat Z=(1-\alpha)c(X)+\alpha\tilde Z\), while the KL regularizes unmixed \(\tilde Z\). Footnote 2 and Appendix D explicitly acknowledge the different laws, yet the p.6 Implication paragraph still says that Vittle maximizes the variational IB lower bound.

\subsection*{Analysis}

Let \(X=\pm2\) equally, label \(Y\) constant, \(c(X)=X\), and \(\tilde Z\mid X\sim\mathcal N(0,1)\) matching its prior independently of \(X\). At \(\alpha=1/2\), the mixed objective is zero for a decoder assigning the sole label probability one. Yet \(\hat Z\mid X=\pm2\sim\mathcal N(\pm1,1/4)\), so \(I(X;\hat Z)>0\), while \(I(Y;\hat Z)=0\). The actual information-bottleneck value is \(-\beta I(X;\hat Z)<0\), so zero is above rather than below it. A finite-softmax likelihood arbitrarily close to zero preserves the strict contradiction.

\subsection*{Finding and scope}

The stated auxiliary-KL objective is not a general lower bound for the representation actually decoded. A useful auxiliary regularizer remains a different interpretation.

\subsection*{Proposed correction}

Use the same mixed law \(\mathcal N((1-\alpha)c+\alpha\mu,\alpha^2\operatorname{diag}\sigma^2)\) in likelihood and KL to an input-independent prior, with any modality-factorization assumptions explicit.

An explicit local correction is described; author adoption and a corresponding published erratum are not established.

\section*{Problem 4: Conditional independence does not eliminate input-shift residuals}

Source: Proposition 3.2 and interpretation after Eq.(7), PDF p.6; Appendix E.2, PDF pp.39–41

Assessment: supported; impact: minor

\subsection*{Source claim (paraphrase)}

After Proposition 3.2, the interpretation drops \(\Delta_{X\mid Z}\) because model output is conditionally independent of input given the representation.

\subsection*{Analysis}

Let domain \(P\) concentrate on input \(a\), and \(Q\) be uniform on \(b,c\). Use one label kernel: a fair label at \(a\), label zero at \(b\), label one at \(c\). Let latents be domain-identical and input-independent, and model output a fair independent label. Conditional independence holds; all retained latent-divergence and marginal entropy-gap terms are zero. True mutual information changes from zero to \(\log2\), so EMID is \(\log2\). The residual is \(2\log2\); retaining it leaves the full proposition uncontradicted. Lowering model entropy can also enlarge an absolute gap to a fixed fair-target entropy.

\subsection*{Finding and scope}

Only the deletion/monotonic interpretation fails; the full printed bound survives this example.

\subsection*{Proposed correction}

Retain the residual unless conditional input laws agree across domains on the relevant representation fibres. Keep entropy terms as absolute gaps rather than assuming lower model entropy always improves them.

An explicit local correction is described; author adoption and a corresponding published erratum are not established.

\section*{Limits}

Preliminary AI assessment; independent human scientific adjudication remains pending.\par

An analytic counterexample or a release discrepancy does not by itself remeasure the paper’s reported experiments.\par

The April 2026 public code postdates the conference; the exact historical checkpoint-training program is not established.\par

\section*{Source identity}

NeurIPS 2025 official proceedings PDF; compared with current arXiv:2505.13946v2

\url{https://papers.nips.cc/paper_files/paper/2025/file/bbc2d0cb01bfafaecebf82d3212ea965-Paper-Conference.pdf}

Reviewed PDF SHA-256: \nolinkurl{f1e9b976f23cabb44e72ccfc1707e1ac59c6ca66c0d46a100675a79e908c4828}

The selected author/publisher source is the same work as the original case, matched by exact title, author identity and official version channel. The exact selected PDF SHA-256 was recomputed in the latest-version pass for Issue \#27. The PDF is a retained author/publisher copy; its original retrieval time is stated only when supported by matching acquisition metadata. File modification time was not substituted.

\end{document}
